\documentclass[12pt]{article}
%---DOCUMENT MARGINS---
\usepackage{geometry} % Required for adjusting page dimensions and margins
\geometry{
	paper=a4paper, % Paper size, change to letterpaper for US letter size
	top=4cm, % Top margin
	bottom=2cm, % Bottom margin
	left=2cm, % Left margin
	right=2cm, % Right margin
	headheight=3cm, % Header height
	%footskip=1.5cm, % Space from the bottom margin to the baseline of the footer
	headsep=0.5cm, % Space from the top margin to the baseline of the header
	%showframe, % Uncomment to show how the type block is set on the page
}
\usepackage{cmap}

\usepackage[T1,T2A]{fontenc}
\usepackage{fontspec}
\setmainfont[
	BoldFont={* Bold},
	ItalicFont={* Italic},
	BoldItalicFont={* BoldItalic}
]{DejaVu Serif Condensed}
\setsansfont{DejaVu Sans}
\setmonofont{Dejavu Sans Mono}
\usepackage[math-style=TeX]{unicode-math}
\setmathfont{Latin Modern Math}

\usepackage[nobottomtitles*]{titlesec}
\titleformat
{\section} % command
{\fontsize{14}{14}\sffamily\bfseries} % format
{} % label
{0pt} % sep
{} % before-code
\titlespacing{\section}{0pt}{0.75em}{0.25em}
\newcommand{\tl}{}
\newcommand{\ml}{}
\newcommand{\problem}[1]{%
	\section[#1]{#1 \fontsize{14}{14}\selectfont{\hfill{\normalfont{
				\emoji{hourglass-not-done}\tl \space\space \emoji{floppy-disk} \ml}}}}
}
\AddToHook{cmd/section/before}{\clearpage}

\usepackage[dvipsnames]{xcolor}
\titleformat
{\subsection} % command
{\fontsize{14}{14}\sffamily\bfseries} % format
{\includesvg[width=0.035\textwidth, inkscapelatex=false]{lion}\space} % label
{0pt} % sep
{} % before-code
\titlespacing{\subsection}{0pt}{0.75em}{0.25em}

\setlength{\parskip}{0.5em}
\setlength{\parindent}{0pt}
\renewcommand{\baselinestretch}{1.2}
\sloppy

\usepackage{listings}
\definecolor{lstbg}{RGB}{250,250,252}
\definecolor{lstframe}{RGB}{220,220,225}
\lstdefinestyle{mystyle}{
	backgroundcolor=\color{lstbg},
	frame=single,
	rulecolor=\color{lstframe},
	basicstyle=\ttfamily\small,
	keywordstyle=\bfseries,
	breaklines=true,
	aboveskip=0.5\baselineskip,
	belowskip=0\baselineskip  
}
\lstset{style=mystyle, language=C++}

\usepackage{fancyhdr}
\pagestyle{fancy}
\setlength{\headwidth}{\textwidth}
\newcommand{\round}{}
\newcommand{\problemName}{}
\newcommand{\lang}{English}
\fancyhead[L]{
	\begin{minipage}{0.4\textwidth}
		\bf{
			EJOI 2025 \round\\
			\problemName\\
			\emoji{globe-with-meridians} \lang
		}
	\end{minipage}
	\vspace{0.5cm}
}
\fancyhead[R]{%
	\begin{minipage}{0.6\textwidth}
		\includesvg[width=\textwidth, inkscapelatex=false]{logo}
	\end{minipage}
	\vspace{0.5cm}
}
\usepackage{lastpage}
\fancyfoot[C]{\problemName\space(\lang) \thepage\ / \pageref*{LastPage}}

\raggedbottom

\usepackage{amsmath}
\usepackage{stmaryrd}

\usepackage{graphicx}
\graphicspath{{./}}
\usepackage[export]{adjustbox}
\usepackage{wrapfig}
\makeatletter
\patchcmd\WF@putfigmaybe{\lower\intextsep}{}{}{\fail}
\AddToHook{env/wrapfigure/begin}{\setlength{\intextsep}{0pt}}
\makeatother
\usepackage[inkscapearea=page, inkscapepath=./svg-inkscape]{svg}
\svgpath{{./}}

\usepackage{makecell}
\usepackage{tabularray}
\AtBeginEnvironment{table}{\vspace{-0.2cm}}
\AtEndEnvironment{table}{\vspace{-0.2cm}}
\usepackage{float}
\usepackage{placeins}
\usepackage{caption}
\captionsetup[table]{
	skip=0.25em,
	singlelinecheck=false, justification=justified
}

\usepackage{enumitem}
\setlist{itemsep=-0.4em, leftmargin=22pt, topsep=-\parskip}
\AfterEndEnvironment{itemize}{\vspace{\parskip}}
\newcommand{\tabitem}{\indent~~\llap{\textbullet}~~}

\usepackage{hyperref}
\hypersetup{
	colorlinks=true,
	citecolor=blue,
	linkcolor=blue,
	urlcolor=cyan,
}

\usepackage{emoji}

\usepackage[english]{babel}

\begin{document}

\renewcommand{\round}{Day 2}
\renewcommand{\problemName}{Task Navigation}
\renewcommand{\lang}{Russian}

\renewcommand{\tl}{$3$ sec.}
\renewcommand{\ml}{$512$ MB}
\problem{\problemName}

Существует \textbf{связный неориентированный простой кактус}\textsuperscript{1} с $N \leq 1000$ вершинами и $M$ ребрами. Его вершины имеют цвета (обозначенные неотрицательными целыми числами от $0$ до $1499$). Изначально все вершины имеют цвет $0$. \textbf{Детерминированный робот без памяти}\textsuperscript{2} исследует граф, перемещаясь от вершины к вершине. Он должен посетить все вершины по крайней мере один раз, а затем завершить работу.

Робот начинает с какой то вершины, которая может быть любой из графа. На каждом шаге он видит цвет своей текущей вершины  и цвета всех соседних верщин \textbf{в некотором фиксированном для текущей вершины порядке} (т. е. повторное посещение вершины даст роботу ту же последовательность соседних вершин, даже если их цвета отличаются от прежних). Робот выполняет одно из следующих двух действий:
\begin{enumerate}
    \item Завершает работу.
    \item Выбирает новый (или, возможно, тот же) цвет для текущей вершины и соседнию вершины, к которой необходимо перейти. Соседняя вершина выбирается индексом от $0$ до $D-1$, где $D$ — количество соседних вершин.
\end{enumerate}

Во втором случае текущая вершина перекрашивается (или, возможно, остается того же цвета), и робот переходит к выбранной соседней вершине. Это повторяется до тех пор, пока робот не завершит работу или не достигнет предела итераций. Робот выигрывает, если он посещает все вершины и затем завершает работу в пределах лимита итераций $L = 3000$ шагов (в противном случае он проигрывает).

Вы должны разработать стратегию для робота, которая сможет решить задачу на любого кактуса. Кроме того, вы должны постараться минимизировать количество различных цветов, используемых в вашем решении. Здесь цвет 0 всегда считается использованным.

\textsuperscript{1}\textit{Связный неориентированный простой кактус — это связный неориентированный простой граф (каждая вершина достижима из любой другой вершины; ребра двунаправлены; нет петлей и кратных ребер), в котором каждое ребро принадлежит не более чем одному простому циклу (простой цикл — это цикл, в котором каждая вершина встречается не более одного раза). Пример показан на рисунке ниже.}

\begin{adjustbox}{center}
	\includesvg[inkscapelatex=false, width=.5\linewidth]{graph}
\end{adjustbox}

\textsuperscript{2}\textit{Робот называется детерминированным и без памяти, если его действия зависят только от текущих входных данных (т.е. он не сохраняет данные между итерациями) и он всегда выбирает одно и то же действие при одинаковых входных данных.}

\subsection{Implementation details}
Стратегия робота должна быть реализована в виде следующей функции:
\begin{lstlisting}
 std::pair<int, int> navigate(int currColor, std::vector<int> adjColors)
\end{lstlisting}

Она принимает в качестве аргументов цвет текущей вершины и цвета всех соседних вершин (в фиксированном порядке). Она должна возвращать пару, первый элемент которой является новым цветом для текущей вершины, а второй элемент — индексом соседней вершины, к которой должен перейти робот. Если же робот должен завершить работу, функция должна возвращать пару $(-1, -1)$.

Эта функция будет вызываться повторно для выбора следующих действий робота. Поскольку она детерминирована, если \texttt{navigate} уже вызывалась с некоторыми аргументами, она никогда не будет вызвана с теми же аргументами снова; вместо этого будет использовано значение, которое оно уже вернула до этого . Кроме того, каждый тест может содержать $T \leq 5$ подтестов (различные графы и/или начальные позиции), и они могут выполняться одновременно (т. е. ваша программа может получать поочередные вызовы по разным подтестам). Наконец, вызовы \texttt{navigate} могут происходить в \textbf{отдельных выполнениях} вашей программы (но иногда они могут происходить и в одном и том же выполнении). Общее количество выполнений вашей программы составляет $P = 100$. Это сделано для того чтобы ваша программа не могла передовать информацию между разными вызовами функции. 

\subsection{Constraints}

\begin{itemize}
\item $3 \leq N \leq 1000$
\item $0 \leq \mathrm{цвет} < 1500$
\item $L = 3000$
\item $T \leq 5$
\item $P = 100$
\end{itemize}

\subsection{Scoring}

Часть $S$ баллов за подзадачу, которую вы получаете, зависит от $C$ — максимального числа различных цветов, используемых в вашем решении (включая цвет $0$) в любом тесте этой подзадачи или любой другой требуемой подзадачи:

\begin{itemize}
\item Если ваше решение не проходит какой-либо подтест, то $S = 0$.
\item Если $C \leq 4$, то $S = 1,0$.
\item Если $4 < C \leq 8$, то $S = 1,0 - 0,6 \frac{C - 4}{4}$.
\item Если $8 < C \leq 21$, то $S = 0,4 \frac{8}{C}$.
\item Если $C > 21$, то $S = 0,15$.
\end{itemize}

\subsection{Subtasks}

\begin{table}[H]
\begin{tblr}{|Q[c,m]|Q[c,m]|Q[c,m]|Q[c,m]|X[c,m]|}
\hline
\textbf{Подзадача} & \textbf{Баллы} & \textbf{\makecell{Требуемые\\подзадачи}} & $N$ & \textbf{Дополнительные\\ограничения}\\
		\hline
$0$ & $0$ & $-$ & $\leq 300$ & Примеры. \\
\hline
$1$ & $6$ & $-$ & $\leq 300$ & Граф является циклом.\textsuperscript{1} \\
\hline
$2$ & $7$ & $-$ & $\leq 300$ & Граф является звездой.\textsuperscript{2} \\ 
\hline
$3$ & $9$ & $-$ & $\leq 300$ & Граф является путем.\textsuperscript{3} \\ 
\hline
$4$ & $16$ & $2 - 3$ & $\leq 300$ & Граф является деревом.\textsuperscript{4} \\
\hline
$5$ & $27$ & $-$ & $\leq 300$ & Все вершины имеют не более $3$ соседних вершин, а вершина, с которой начинает робот, имеет ровно $1$ соседнюю вершину.\\
\hline
$6$ & $28$ & $0 - 5$ & $\leq 300$ & $-$ \\
\hline
$7$ & $7$ & $0 - 6$ & $-$ & $-$ \\
\hline
\end{tblr}
\caption*{\textsuperscript{1}Граф является циклом если имеет ребра: $\left(i, (i+1) \bmod N\right)$ для  $0 \leq i < N$. \\
\textsuperscript{2}Граф является звездой если имеет ребра: $\left(0, i\right)$ для  $1 \leq i < N$. \\
\textsuperscript{3}Граф является путем если имеет ребра: $\left(i, i+1\right)$ для $0 \leq i < N - 1$. \\
\textsuperscript{4}Дерево — это связный граф без циклов.}
\end{table}
\FloatBarrier

\subsection{Example}

Рассмотрим пример графа из изображения в условии, который имеет $N=7$, $M=8$ и ребра $(0, 1)$, $(1, 2)$, $(2, 0)$, $(2, 3)$, $(3, 4)$, $(4, 2)$, $(3, 5)$ и $(2, 6)$. Кроме того, поскольку порядок соседних вершин иммет значение, мы приводим его в этой таблице:


\begin{table}[H]
\begin{tblr}{|Q[c,m]|Q[c,m]|}
\hline
\textbf{\makecell{Node}} & \textbf{Adjacent nodes} \\
\hline
$0$ & $2, 1$ \\
\hline
$1$ & $2, 0$ \\
\hline
$2$ & $0, 3, 4, 6, 1$ \\
\hline
$3$ & $4, 5, 2$ \\
\hline
$4$ & $2, 3$ \\
\hline
$5$ & $3$ \\
\hline
$6$ & $2$ \\
\hline
\end{tblr}
\end{table} 



Предположим, что робот начинает с вершины $5$. Тогда ниже приведена одна из возможных (неудачных) последовательностей вызовов:

\begin{table}[H]
\begin{tblr}{|Q[c,m]|Q[c,m]|Q[c,m]|X[c,m]|Q[c,m]|}
\hline
\textbf{\makecell{\#}} & \textbf{\makecell{Цветы}} & \textbf{\makecell{Текущая вершина}} & \textbf{\makecell{Вызов \texttt{navigate}}} & \textbf{Ответ функции} \\
\hline
$1$ & $0, 0, 0, 0, 0, 0, 0$ & $5$ & \texttt{navigate(0, \{0\})} & \texttt{\{1, 0\}} \\
\hline
$2$ & $0, 0, 0, 0, 0, 1, 0$ & $3$ & \texttt{navigate(0, \{0, 1, 0\})} & \texttt{\{4, 2\}} \\
\hline
$3$ & $0, 0, 0, 4, 0, 1, 0$ & $2$ & \texttt{navigate(0, \{0, 4, 0, 0, 0\})} & \texttt{\{0, 3\}} \\
\hline
$4$ & $0, 0, 0, 4, 0, 1, 0$ & $6$ & \textsuperscript{1}\texttt{navigate(0, \{0\})} & \texttt{\{1, 0\}} \\
\hline
$5$ & $0, 0, 0, 4, 0, 1, 1$ & $2$ & \texttt{navigate(0, \{0, 4, 0, 1, 0\})} & \texttt{\{8, 0\}} \\
\hline
$6$ & $0, 0, 8, 4, 0, 1, 1$ & $0$ & \texttt{navigate(0, \{8, 0\})} & \texttt{\{3, 0\}} \\
\hline
$7$ & $3, 0, 8, 4, 0, 1, 1$ & $2$ & \texttt{navigate(8, \{3, 4, 0, 1, 0\})} & \texttt{\{2, 2\}} \\
\hline
$8$ & $3, 0, 2, 4, 0, 1, 1$ & $4$ & \texttt{navigate(0, \{2, 4\})} & \texttt{\{1, 1\}} \\
\hline
$9$ & $3, 0, 2, 4, 1, 1, 1$ & $3$ & \texttt{navigate(4, \{1, 1, 2\})} & \texttt{\{-1, -1\}} \\
\hline
\end{tblr}
\end{table}

Здесь робот использовал в общей сложности $6$ различных цветов: $0$, $1$, $2$, $3$, $4$ и $8$ (обратите внимание, что $0$ считался бы использованным, даже если бы робот никогда не возвращал цвет $0$, поскольку все узлы начинают с цвета $0$). Робот проработал $9$ итераций, прежде чем завершить работу. Однако он провалился, поскольку завершил работу, не посетив узел $1$.

\textsuperscript{1}Обратите внимание, что вызов \texttt{navigate} на итерации $4$ на самом деле не произойдет. Это потому, что он эквивалентен вызову на итерации $1$, поэтому грейдер просто повторно использует возвращаемое значение вашей функции из этого вызова. Однако это все равно считается итерацией робота.

\subsection{Sample grader}

Грейдер не запускает несколько запусков вашей программы, поэтому все вызовы \texttt{navigate} будут находиться в одном запуске вашей программы.
\begin{itemize}
\item строка $1$: два целых числа — $N$ и $M$; 
\item строка $2+i$ (для $0 \leq i < M$): два целых числа — $A_i$ и $B_i$, вершины соединяемыми ребром $i$ ($0 \leq A_i, B_i < N$). 
\end{itemize}

Затем грейдер выведет количество различных цветов, использованных в вашем решении, и количество итераций, необходимых для его завершения. В противном случае, если ваше решение не прошло проверку, будет выведено сообщение об ошибке.

По умолчанию, грейдер выводит подробную информацию о том, что робот видит и делает на каждой итерации. Вы можете отключить эту функцию, изменив значение \texttt{DEBUG} с \texttt{true} на \texttt{false}.

\end{document}
